% !TeX program = pdflatex
%
% The primitive solutions of x^4 + 2y^3 + z^3 = 0.
% Companion to the repository https://github.com/marchfederico/fermat433
%
\documentclass[11pt,reqno]{amsart}

\usepackage[OT2,T1]{fontenc}   % OT2 supplies the Cyrillic "Sha" for \Sha below
\usepackage[utf8]{inputenc}
\usepackage{amsmath,amssymb,amsthm}
\usepackage{mathtools}
\usepackage{array}
\usepackage{booktabs}
\usepackage{enumitem}
\usepackage[margin=1.15in]{geometry}
\usepackage[colorlinks=true,linkcolor=blue,citecolor=blue,urlcolor=blue]{hyperref}

% ---------------------------------------------------------------- theorems ---
\theoremstyle{plain}
\newtheorem{theorem}{Theorem}[section]
\newtheorem{proposition}[theorem]{Proposition}
\newtheorem{lemma}[theorem]{Lemma}
\newtheorem{corollary}[theorem]{Corollary}
\theoremstyle{definition}
\newtheorem{remark}[theorem]{Remark}
\newtheorem{convention}[theorem]{Convention}

% ----------------------------------------------------------------- notation ---
\newcommand{\Q}{\mathbb{Q}}
\newcommand{\Z}{\mathbb{Z}}
\newcommand{\R}{\mathbb{R}}
\newcommand{\Co}{\mathbb{C}}
\newcommand{\Ad}{\mathbb{A}}
\newcommand{\PP}{\mathbb{P}}
\newcommand{\FF}{\mathbb{F}}
\newcommand{\Qbar}{\overline{\Q}}
\newcommand{\Ocal}{\mathcal{O}}
\newcommand{\pp}{\mathfrak{p}}
\newcommand{\qq}{\mathfrak{q}}
\newcommand{\PPr}{\mathfrak{P}}
\newcommand{\aaa}{\mathfrak{a}}

\DeclareMathOperator{\NS}{NS}
\DeclareMathOperator{\rk}{rk}
\DeclareMathOperator{\Aut}{Aut}
\DeclareMathOperator{\End}{End}
\DeclareMathOperator{\Jac}{Jac}
\DeclareMathOperator{\Sel}{Sel}
\DeclareMathOperator{\disc}{disc}
\DeclareMathOperator{\Gal}{Gal}
\DeclareMathOperator{\Ind}{Ind}
\DeclareMathOperator{\sgn}{sgn}
\DeclareMathOperator{\stdrep}{std}
\DeclareMathOperator{\Log}{Log}
\DeclareMathOperator{\Nm}{N}
\DeclareMathOperator{\cont}{cont}
\DeclareMathOperator{\Sym}{S}
\newcommand{\tors}{\mathrm{tors}}
\newcommand{\fake}{\mathrm{fake}}

% The Tate--Shafarevich group.  \Sha below uses the wncyr fonts, which live in
% TeX Live's "cyrillic" collection (present on Overleaf and in any full
% installation).  If they are missing, comment out the two \Declare... lines
% and uncomment the \newcommand fallback instead.
\DeclareSymbolFont{cyrletters}{OT2}{wncyr}{m}{n}
\DeclareMathSymbol{\Sha}{\mathord}{cyrletters}{"58}
%\newcommand{\Sha}{\mathrm{Sha}}

\numberwithin{equation}{section}

% Long \texttt{} script names cannot hyphenate; give TeX room to stretch.
\emergencystretch=3em

% ------------------------------------------------------------------- front ---
\title[Quadratic Chabauty through a hidden involution]
      {Quadratic Chabauty through a hidden involution:\\
       the primitive solutions of $x^{4}+2y^{3}+z^{3}=0$}

\author{Marcello Federico}
\email{marcellof@gmail.com}
\date{September 2026}

\subjclass[2020]{Primary 11D41; Secondary 11G30, 11G50, 11Y50, 14G05, 14H45}
\keywords{Generalised Fermat equation, quadratic Chabauty, bielliptic genus two
curve, N\'eron--Severi group, $p$-adic height, Mordell--Weil sieve, Bolza curve}

\begin{document}

\begin{abstract}
We determine all primitive integer solutions of $x^{4}+2y^{3}+z^{3}=0$: they are
$(\pm1,0,-1)$, $(\pm1,-1,1)$ and $(\pm43,-203,237)$. This was the smallest
generalised Fermat equation left uninvestigated ($H=40$) in the survey of
Ratcliffe and Grechuk. A descent over $\Q(\sqrt[3]{2})$ reduces the equation to
rational points on six genus-$2$ curves, four of which are settled by two-cover
descent and by a classical Chabauty argument at an inert prime. The remaining
pair, $B^{-}\colon y^{2}=-C_{3}(x)$ and $C^{+}\colon y^{2}=2C_{3}(x)$ with
$C_{3}(x)=17x^{6}-72x^{5}+90x^{4}+40x^{3}-180x^{2}+144x-40$, have Jacobians of
Mordell--Weil rank $2$, equal to the genus, and N\'eron--Severi rank $1$ over
$\Q$ and over $\Q(\sqrt{-2})$; neither classical Chabauty nor quadratic
Chabauty over those fields applies.

We show that these curves are $\Q$-twists of the Bolza curve and carry an
involution $\tau$ defined over a field $L$ of degree $12$, whose N\'eron--Severi
class is the unique class on which complex conjugation acts by $-1$ --- exactly
the class counted by the term $\rk\NS(J_{\Qbar})^{c=-1}$ in the finiteness
criterion of Balakrishnan--Dogra. The bielliptic quadratic Chabauty function of
Balakrishnan--Besser--Bianchi--M\"uller over $L$, formed with the unique idele
class character lying in the matching Galois-isotypic component, restricts to a
locally analytic function on $X(\Q_{p})$ that takes a single value on $X(\Q)$.
Computing its zeros at $p=499$, and independently at $p=1459$, and applying a
Mordell--Weil sieve leaves exactly the four known rational points on each curve.

\medskip
\noindent\emph{Provenance.} This work was carried out in a human--AI
collaboration with Claude (Anthropic): the design of the computations, the
scripts and this write-up were all produced in that collaboration. Every
assertion below is backed by a script and a log in the public repository, and
the record of how the argument developed, including the errors found and
corrected along the way, is kept there as well. The result has not been
refereed. Credits and status in full: \S\ref{subsec:status}.
\end{abstract}

\maketitle

% ==============================================================================
\section{Introduction}
% ==============================================================================

\subsection{The result}

\begin{theorem}\label{thm:main}
The integer solutions of
\begin{equation}\label{eq:main}
  x^{4}+2y^{3}+z^{3}=0, \qquad \gcd(x,y,z)=1,
\end{equation}
are exactly
\[
  (x,y,z)=(\pm1,\,0,\,-1),\qquad (\pm1,\,-1,\,1),\qquad (\pm43,\,-203,\,237).
\]
\end{theorem}

Equation \eqref{eq:main} is a generalised Fermat equation of signature
$(4,3,3)$ and height $H=40$. In the survey of Ratcliffe and Grechuk
\cite{RatcliffeGrechuk} it is the smallest equation recorded as
uninvestigated. Villines \cite{Villines} reduced it, with Magma, to rational
points on six curves of genus $2$ and settled four of them, leaving open the
two whose Jacobians have Mordell--Weil rank equal to the genus. The present
paper closes that gap, and re-proves the remaining parts independently; all
computations here run on PARI/GP alone.

The proof has four parts.
\begin{enumerate}[label=(\Alph*),leftmargin=2.4em]
\item A descent over $K=\Q(\sqrt[3]{2})$ reducing \eqref{eq:main} to rational
      points on six genus-$2$ curves (\S\ref{sec:descent}).
\item The four accessible curves: three have empty fake $2$-Selmer set, and the
      fourth has rank $1$ and is settled by Chabauty at an inert prime together
      with a Mordell--Weil sieve (\S\ref{sec:easy}).
\item The hard pair $B^{-}$, $C^{+}$, where $\rk J(\Q)=2=g$, by quadratic
      Chabauty (\S\S\ref{sec:structure}--\ref{sec:sieve}). This is the new
      contribution.
\item The rank bounds used in (C), by a $2$-descent over $\Q$
      (\S\ref{sec:rank}).
\end{enumerate}

\subsection{The obstruction, and the idea}

Write $X$ for either curve of the hard pair and $J=\Jac(X)$. Classical
Chabauty--Coleman needs $\rk J(\Q)<g$, and here $\rk J(\Q)=2=g$. Quadratic
Chabauty in the form of Balakrishnan--Dogra \cite{BDii} needs a supply of
N\'eron--Severi classes: the depth-$2$ Chabauty--Kim set $X(\Q_{p})_{2}$ is
finite as soon as
\begin{equation}\label{eq:BDcount}
  \rk J(\Q) \;<\; g-1+\rk\NS(J)+\rk\NS(J_{\Qbar})^{c=-1},
\end{equation}
$c$ denoting complex conjugation. We show (Proposition~\ref{prop:NS}) that
$\NS(J_{\Qbar})\otimes\Q$ has rank $4$ and, as a representation of
$\Gal(M/\Q)\cong \Sym_{4}\times\Z/2$ for an explicit field $M$ of degree $48$,
is
\[
  \NS(J_{\Qbar})\otimes\Q \;\cong\; \mathbf{1}\oplus(\sgn_{F}\otimes\sgn\otimes\stdrep),
\]
so that $\rk\NS(J)=1$ and complex conjugation has eigenvalues $(1,1,1,-1)$. The
right-hand side of \eqref{eq:BDcount} is therefore $2-1+1+1=3>2$: finiteness
holds, but the only usable class is the single anti-invariant one, and it is
not defined over $\Q$. The same computation shows $\rho(J/F)=1$ for
$F=\Q(\sqrt{-2})$, so quadratic Chabauty over $F$ --- the route suggested in
\cite{Villines} --- has no class to use either.

The anti-invariant class turns out to be visible. The M\"obius transformations
permuting the six roots of $C_{3}$ form $\Sym_{4}$, so $\Aut(X_{\Qbar})\cong
\mathrm{GL}_{2}(\FF_{3})$ has order $48$ and $X$ is a $\Q$-twist of the Bolza
curve $y^{2}=x^{5}-x$. Six of the nine involutory M\"obius maps lift to genuine
involutions $\tau$ of $X$, and these are defined over a field
\[
  L=\Q\bigl(\sqrt[3]{2},\;\sqrt{\sqrt[3]{2}-1},\;\sqrt{\kappa}\bigr),
  \qquad [L:\Q]=12,
\]
whose degree is exactly what the N\'eron--Severi character predicts (the
stabiliser of the anti-invariant class has order $4$ in a group of order $48$).
Over $L$ the curve is bielliptic, with quotients $E_{1},E_{2}$ of $j$-invariant
$8000$.

This puts us in the situation of the explicit bielliptic quadratic Chabauty of
Balakrishnan--Besser--Bianchi--M\"uller \cite{BBBM}, whose function is built
from an idele class character $\chi$ of the base field. The point is a
\emph{selection rule}: for $z\in X(\Q)\subset X(L)$ all Galois-conjugate places
contribute the same local terms, so the identity survives on $\Q$-points only
when the class and the character lie in the same Galois-isotypic component. The
cyclotomic character lies in the trivial component and yields $0=0$. In the
$3$-dimensional component $\sgn_{F}\otimes\sgn\otimes\stdrep$ that matches
$\tau_{*}$ there is, we show, exactly one character $\chi_{Z}$ up to scaling
(Proposition~\ref{prop:character}); and it vanishes on the uniformisers of every
prime of $L$ above $2$ and $3$, which is what makes the resulting identity
independent of the reduction types at the bad primes.

The outcome (Theorem~\ref{thm:identity}) is that
\[
  \rho(z)\;=\;\sum_{i=1}^{12}c_{i}\bigl[\lambda_{i}(\varphi_{1}z)-\lambda_{i}(\varphi_{2}z)\bigr]
  \;-\;2\sum_{i=1}^{12}c_{i}\log_{p}\sigma_{i}(X(z))\;-\;a(z)^{\mathsf T}Ga(z)
\]
is a locally analytic function on $X(\Q_{p})$ taking the single value
$-\Omega_{397}$ on $X(\Q)$. Its zero set is computed disk by disk, and a
Mordell--Weil sieve cuts the survivors down to the four known rational points on
each curve.

\subsection{What is new}

\begin{enumerate}[label=(\arabic*),leftmargin=2.4em]
\item An explicit quadratic Chabauty computation for a curve over $\Q$ whose
      extra N\'eron--Severi class is defined only over a field of degree $12$,
      found from the automorphism group of a twist of the Bolza curve, with
      the field of definition matching the Galois character of
      $\NS(J_{\Qbar})$.
\item The identification of the character by linear algebra from the action of
      the involutions on holomorphic forms together with the $p$-adic
      logarithms of the units of $L$, with the vanishing above $2$ and $3$ as a
      by-product.
\item $p$-adic heights at twelve places of a degree-$12$ field, computed through
      integer models congruent to the completions modulo $p^{30}$, and
      division-polynomial recurrences for $\psi_{514}$ in place of a polynomial
      of degree ${\approx}130{,}000$.
\item An index-free Mordell--Weil sieve: $p$-integrality of the coefficient
      vector removes the saturation hypothesis usually needed.
\item The observation that the endgame proposed in \cite{Villines} over
      $\Q(\sqrt{-2})$ cannot work, since $\rho(J/\Q(\sqrt{-2}))=1$.
\end{enumerate}

\subsection{What the proof rests on}\label{subsec:rests}

Parts (A), (B) and (D) are self-contained exact or $p$-adic computations. Part
(C) additionally imports:

\begin{enumerate}[label=(\roman*),leftmargin=2.4em]
\item $p$-adic heights on elliptic curves over a number field for a
      non-cyclotomic idele class character: that $h^{\chi}$ is a quadratic form,
      that it decomposes into local terms, that the local term at $w\nmid p$ of
      good reduction is $\chi_{w}(\pi_{w})\max(0,-v_{w}(x(Q)))$ on a model
      minimal at $w$, and that at $\pp\mid p$ of good ordinary reduction it is
      the canonical Mazur--Tate height computed from the $\sigma$-function
      \cite{MazurTate,Bianchi,BBBM}.
\item The bielliptic quadratic Chabauty identity of \cite[Thm.~1.5]{BBBM} over a
      number field, with the normalisation of the logarithmic term as in
      \cite[Thm.~2.3]{BianchiPadurariu}.
\item The finiteness theorem of \cite[Prop.~2]{BDii} (quoted in
      \cite{Dogra}); this is used only as motivation, the finiteness of the
      zero set being established computationally by Strassman bounds.
\item Standard Coleman integration on residue disks.
\item PARI's rigorous $p$-adic precision tracking. The only truncation is the
      series tail on each residue disk, controlled by the observed linear growth
      of coefficient valuations and confirmed at higher precision.
\end{enumerate}

\subsection{Provenance and credits}\label{subsec:status}

This work was done in collaboration with Claude (Anthropic): the design of the
computations, the PARI/GP and Python scripts, and the present write-up were all
produced in that collaboration. The author is nonetheless answerable for
everything asserted here, which is what the name on the title page means.
Claude is not an author and could not be one: it can neither stand behind the
result nor answer for it if the result is wrong.

The reader should treat the mathematics accordingly --- that is,
critically, and with the scripts and logs in hand. Everything needed to do so is
public (\S\ref{sec:code}): every proposition names the script and the log that
establish it (Table~\ref{tab:scripts}), and \texttt{RESULTS.md} in the
repository is a chronological record of how the argument developed, including
the several errors that were found and corrected and the checks that caught
them.

This is therefore a research record, not a refereed result. The computations
have been repeated at a second prime and cross-checked against the independent
Magma computations of \cite{Villines} wherever the two overlap, and they agree
throughout; but the argument of \S\ref{sec:identity}, and in particular the
precise relationship between the identity proved there and
\cite[Thm.~1.5]{BBBM}, has not been examined by anyone else. The questions the
author would most like an expert to consider are collected in
\texttt{REFEREE\_NOTES.md} in the repository.

\begin{convention}
Throughout, $\iota$ denotes the hyperelliptic involution $(x,y)\mapsto(x,-y)$;
$\theta=\sqrt[3]{2}$; $K=\Q(\theta)$; $F=\Q(\sqrt{-2})$; and
\[
  C_{3}(x)=17x^{6}-72x^{5}+90x^{4}+40x^{3}-180x^{2}+144x-40 .
\]
For a curve $C\colon y^{2}=f(x)$ of genus $2$ we write $J=\Jac(C)$. Every
assertion labelled with a script name is established by that script; the
scripts and their logs are in the repository (\S\ref{sec:code}), and
Table~\ref{tab:scripts} indexes them.
\end{convention}

% ==============================================================================
\section{Descent to six curves of genus two}\label{sec:descent}
% ==============================================================================

Let $\theta=\sqrt[3]{2}$ and $K=\Q(\theta)$. Then $K$ has class number $1$
(certified with \texttt{bnfcertify}), fundamental unit $\varepsilon=\theta-1$ of
norm $1$, and $2$ and $3$ are totally ramified.

\begin{lemma}\label{lem:A1}
Let $(x,y,z)$ be a primitive solution of \eqref{eq:main}. Then $\gcd(y,z)=1$ and
$\alpha:=z+y\theta$ satisfies $\Nm_{K/\Q}(\alpha)=z^{3}+2y^{3}=-x^{4}$. The
three conjugates of $\alpha$ are pairwise coprime away from the primes above
$2$ and $3$, and primitivity forces the common factors there to be trivial;
hence $(\alpha)=\aaa^{4}$ for an ideal $\aaa$. Since $h_{K}=1$ and
$\Ocal_{K}^{\times}=\langle-1,\varepsilon\rangle$ we may write
\[
  \alpha=-\varepsilon^{k}\gamma^{4},\qquad \gamma\in\Z[\theta],\quad k\in\{0,1,2,3\},
\]
the sign being normalised by replacing $x$ with $-x$ if necessary.
\end{lemma}

\begin{proof}
The ideal argument is the standard one for norm forms of cubic fields; see
\texttt{descent\_explore.py}.
\end{proof}

\begin{lemma}\label{lem:A2}
Write $\gamma=a+b\theta+c\theta^{2}$. The coefficient of $\theta^{2}$ in
$-\varepsilon^{k}\gamma^{4}$ must vanish, which is a smooth plane quartic
$C_{k}(a,b,c)=0$ over $\Q$, of genus $3$. Primitivity forces $\gamma$ to be a
unit at the primes above $2$ and $3$, that is, $a$ odd and
$a+2b+c\not\equiv0\pmod 3$. The curve $C_{3}$ has no $2$-adic points, and
$C_{2}$ has no $2$-adic point with $a$ odd. Hence every primitive solution comes
from a rational point of $C_{0}$ or $C_{1}$ satisfying the unit conditions, with
$x=\Nm(\gamma)$ and $z+y\theta=-\varepsilon^{k}\gamma^{4}$.
\end{lemma}

\begin{proof}
\texttt{descent\_explore.py}, \texttt{diagonal.py}; the explicit quartics are
recorded in \texttt{RESULTS.md}.
\end{proof}

\begin{lemma}[\'Etale double covers and the twist set]\label{lem:A3}
Put $\delta=\gamma^{2}$. The condition $[\theta^{2}](-\varepsilon^{k}\delta^{2})=0$
is a smooth conic $Q_{k}(d_{0},d_{1},d_{2})=0$ in the coordinates of
$\delta=d_{0}+d_{1}\theta+d_{2}\theta^{2}$ (Gram determinant $1$ for both $k$),
with the rational point $(1:0:0)$. Parametrising by the lines through that point,
\begin{align*}
 k=0:\quad & Q_{0}=-d_{1}^{2}-2d_{0}d_{2},\qquad
             \delta_{0}(s,t)=(s^{2},\,-2st,\,-2t^{2}),\\
           & f_{0}(s,t)=\Nm(\delta_{0})=s^{6}-40s^{3}t^{3}-32t^{6};\\[3pt]
 k=1:\quad & Q_{1}=d_{1}^{2}-2d_{0}d_{1}+2d_{0}d_{2}-2d_{2}^{2},\qquad
             \delta_{1}(s,t)=(2t^{2}-s^{2},\,2s(t-s),\,2t(t-s)),\\
           & f_{1}(s,t)=-t^{6}C_{3}(s/t).
\end{align*}
Every rational point of the conic is $\delta_{k}(s:t)$ for a unique
$(s:t)\in\PP^{1}(\Q)$; with $s,t$ coprime integers,
$\delta=\mu\cdot\delta_{k}(s,t)$ for a unique $\mu\in\Q^{\times}$, and taking
norms, $\mu f_{k}(s,t)=(\Nm(\gamma)/\mu)^{2}$. Then $\mu\in\{\pm1,\pm\tfrac12\}$,
and $\mu=\pm1$ when $k=0$. Consequently every primitive solution gives a
rational point on one of the six curves
\begin{align}
 &A^{\pm}\colon y^{2}=\pm\bigl(x^{6}-40x^{3}-32\bigr), \label{eq:Apm}\\
 &B^{\mp}\colon y^{2}=\mp C_{3}(x), \qquad
  C^{\mp}\colon y^{2}=\mp 2\,C_{3}(x). \label{eq:BCpm}
\end{align}
\end{lemma}

\begin{proof}
Each point of $C_{k}(\Q)$ gives the rational point
$\bigl(s/t,\ \Nm(\gamma)/(\mu t^{3})\bigr)$ on $H_{\mu}\colon w^{2}=\mu f_{k}(x,1)$,
and the point of $C_{k}$ is recovered from it by $\delta=\mu\delta_{k}(s,t)$,
$\gamma=\sqrt{\delta}$. It remains to pin down the twist class of $\mu$.

\emph{(i) Content.} For coprime $(s,t)$ the content of $\delta_{k}(s,t)$ is $1$
if $s$ is odd and $2$ if $s$ is even. Indeed an odd prime dividing all three
coordinates would divide both $s$ and $t$: for $k=0$ it divides $s^{2}$ and
$2t^{2}$; for $k=1$ it divides $2s(t-s)$ and $2t(t-s)$, hence $t-s$, hence
$2t^{2}-s^{2}\equiv t^{2}$. The claim modulo $4$ follows by inspecting the
coordinates $(s^{2},2st,2t^{2})$, resp.\ $(2t^{2}-s^{2},\dots)$ with the other
two even.

\emph{(ii) $\gamma^{2}$ is primitive.} $\gamma$ is primitive; if $p$ divided all
coordinates of $\gamma^{2}$ then $\gamma$ would be nilpotent in
$\Z[\theta]/(p)$, which is reduced for $p\nmid6$, while for $p=2,3$ the
nilpotents are $(\theta)$ and $(\theta+1)$, excluded by the unit conditions of
Lemma~\ref{lem:A2}. Hence
$\mu=\pm1/\cont(\delta_{k}(s,t))\in\{\pm1,\pm\tfrac12\}$.

\emph{(iii) The case $k=0$.} The $\theta^{0}$-coordinate of $\delta=\gamma^{2}$
is $a^{2}+4bc$, which is odd, while that of $\mu\delta_{0}(s,t)$ is $\mu s^{2}$;
so $s$ is odd, the content is $1$, and $\mu=\pm1$. For $k=1$ both contents
occur. Multiplying $w^{2}=\pm\tfrac12 f_{1}$ by $4$ gives the models
\eqref{eq:BCpm}.
\end{proof}

These are the six curves of \cite{Villines}. The known rational points transport
as follows (\texttt{etale.py}): the points $(0,1,0)$, $(1,0,-1)$, $(1,0,0)$ of
$C_{0}$ give $(0:1)$, $(1:-1)$, $(1:0)$ on $A^{+}$; and $(1,-2,0)$, $(1,0,0)$,
$(1,-1,2)$, $(1,1,0)$ of $C_{1}$ give $(-1,\pm15)$, $(1,\pm1)$ on $B^{-}$ and
$(6/5,\pm172/125)$, $(2,\pm12)$ on $C^{+}$. The point $(6/5,\pm172/125)$ is the
solution $(\pm43,-203,237)$.

\begin{remark}
The sign of $\mu$ is the sign of $f_{k}(s,t)$. The point of $C_{0}$ with
$\gamma=\theta$, i.e.\ $\delta=\theta^{2}$ and $(s,t)=(0,1)$, has $s$ even and is
excluded by the unit condition ``$a$ odd''; it corresponds to an imprimitive
solution. So the content-$2$ chart for $k=0$, which would land on the twists
$y^{2}=\pm2(x^{6}-40x^{3}-32)$, never arises from a primitive solution.
\end{remark}

% ==============================================================================
\section{The four accessible curves}\label{sec:easy}
% ==============================================================================

\begin{proposition}\label{prop:B1}
$A^{-}(\Q)=B^{+}(\Q)=C^{-}(\Q)=\emptyset$.
\end{proposition}

\begin{proof}
Let $C\colon y^{2}=f(x)$ be one of the three, with leading coefficient
$f_{6}\in\{-1,17,-34\}$ respectively, let $L_{6}=\Q[x]/(f)$ --- so
$L_{6}\cong\Q[u]/(u^{6}-3u^{4}-2u^{3}+3u^{2}+6u+1)$ for $A^{\pm}$ and
$\Q[u]/(u^{6}-3u^{4}+3u^{2}-3)$ for $B^{\pm},C^{\pm}$, both of class number $1$
(certified by \texttt{bnfcertify}) --- let $\vartheta$ be the image of $x$, and
let $S$ be the set of primes dividing $2\disc(f)f_{6}$, namely $\{2,3\}$
resp.\ $\{2,3,17\}$.

Since $f$ has no rational root and $f_{6}$ is not a square, $C$ has no
Weierstrass and no rational point at infinity. For an affine point
$(x,y)\in C(\Q)$,
\[
  \Nm_{L_{6}/\Q}(x-\vartheta)=f(x)/f_{6}=y^{2}/f_{6}\equiv f_{6}
  \pmod{\Q^{\times2}},
\]
and the class of $x-\vartheta$ in $L_{6}^{\times}/L_{6}^{\times2}\Q^{\times}$ is
unramified outside $S$. (This is \cite[Lemma 4.3]{BruinStoll}: for $p\notin S$,
$f$ is squarefree mod $p$ with unit leading coefficient, so $x\bmod p$ is a root
of exactly one --- linear --- factor of $f$ mod $p$; a single prime $\pp\mid p$
divides $x-\vartheta$, and $v_{\pp}(x-\vartheta)=v_{p}(y^{2}/f_{6})$ is even;
and if $v_{p}(x)<0$ then $(x-\vartheta)/x\equiv1$ modulo every prime above $p$.)
Multiplying by an element of $\Q^{\times}$ changes the norm by a sixth power, so
a rational point would produce an $S$-unit $\delta$ of $L_{6}$ with
$\Nm(\delta)\in f_{6}\Q^{\times2}$.

The norms of the $S$-units modulo squares span the subgroup
$\langle -2,-3\rangle$ of $\Q(S,2)$ for the sextic field of $A^{\pm}$ and
$\langle -2,-3,-17\rangle$ for that of $B^{\pm},C^{\pm}$
(\texttt{sel2/certify.gp}). Neither subgroup contains $-1$ (needed for
$A^{-}$), $17$ (for $B^{+}$) or $-34$ (for $C^{-}$) --- while, as they must,
they do contain $-17$ (for $B^{-}$) and $34$ (for $C^{+}$). Hence
$C(\Q)=\emptyset$.
\end{proof}

This is the emptiness of the fake $2$-Selmer set of Bruin--Stoll
\cite{BruinStoll} at their norm step, and reproduces the corresponding result of
\cite{Villines}. The script \texttt{sel2/fakeselmerset.gp} implements the full
fake Selmer set --- norm condition, real place, and the exact local images
$\mu_{p}(C(\Q_{p}))$ by recursive subdivision of $p$-adic disks --- and returns
$W=\emptyset$ for these three curves already at the norm step.

\begin{proposition}\label{prop:B2}
$A^{+}(\Q)=\{\infty_{\pm},\,(-1,\pm3)\}$.
\end{proposition}

\begin{proof}
\emph{(a) Rank.} $f=x^{6}-40x^{3}-32$ is irreducible, so $J(\Q)[2]=0$, and
\texttt{sel2/pipeline.gp} gives $\dim\Sel^{2}(J/\Q)=1$ --- the sextic field has
the quadratic subfield $\Q(\sqrt3)$, since $f=(x^{3}-20)^{2}-432$, so here
$\Sel^{2}=\Sel^{2}_{\fake}$ --- whence $\rk J(\Q)\le1$. The class
$D=[(-1,3)-\infty_{+}]$ has infinite order, its images on the elliptic quotients
below having infinite order. So $\rk J(\Q)=1$. The torsion subgroup is
$\Z/3$, generated by $T=[\infty_{-}-\infty_{+}]$: the gcd of $\#J(\FF_{\ell})$
over $5\le\ell<400$ is $3$, and $T$ has order $3$.

\emph{(b) Bielliptic structure.} Since $f(c/x)=-32f(x)/x^{6}$ for $c^{3}=-32$,
$X=A^{+}$ carries the involution $\tau(x,y)=(c/x,\,4\sqrt{-2}\,y/x^{3})$ over
$M_{0}=\Q(\sqrt[3]{2},\sqrt{-2})=\Q(\sqrt[6]{-2})$, with fixed points
$x_{\pm}=\pm w^{5}$ where $w^{6}=-2$, $\sqrt[3]{2}=-w^{2}$, $\sqrt{-2}=-w^{3}$
and $c=-2w^{4}$. In the coordinates $\mathcal{X}=(x-x_{+})/(x-x_{-})$,
$\mathcal{Y}=y(1-\mathcal{X})^{3}$ the curve becomes
$\mathcal{Y}^{2}=A\mathcal{X}^{6}+B\mathcal{X}^{4}+C\mathcal{X}^{2}+D$ with
$A=-64+160w^{3}$, $B=-960-480w^{3}$, $C=-960+480w^{3}$, $D=-64-160w^{3}$ (all in
$\Q(\sqrt{-2})$), and $\tau\iota$ is $(\mathcal{X},\mathcal{Y})\mapsto(-\mathcal{X},\mathcal{Y})$.
The quotients
\[
  E_{1}\colon v^{2}=u^{3}+Bu^{2}+ACu+A^{2}D,\qquad
  E_{2}\colon v^{2}=u^{3}+Cu^{2}+BDu+AD^{2},
\]
with $\varphi_{1}=(A\mathcal{X}^{2},A\mathcal{Y})$ and
$\varphi_{2}=(D/\mathcal{X}^{2},D\mathcal{Y}/\mathcal{X}^{3})$, have $j=8000$ (CM
by $\Z[\sqrt{-2}]$), conductor of norm $9$ over $M_{0}$,
$E_{i}(M_{0})_{\tors}\cong\Z/6$, and $E_{2}$ is the quadratic twist of $E_{1}$
by $-2$. $J$ is isogenous to $E_{1}\times E_{2}$ over $M_{0}$
(\texttt{aplus/bi\_Aplus.gp}).

\emph{(c) Chabauty at $p=5$.} The prime $5$ is inert in $\Q(\sqrt{-2})$ and has
three primes of degree $2$ in $M_{0}$; fix $\pp\mid5$, so
$M_{0,\pp}=\Q_{25}=\Q_{5}(\sqrt2)$. Let
$\mathcal{L}_{i}(z)=\Log_{E_{i},\pp}(\varphi_{i}(z))\in\Q_{25}$ be the
formal-group logarithms, computed from the formal group law of the
$a_{1}=a_{3}=0$ model after multiplication by $\#E_{i}(\FF_{25})=36$
(\texttt{aplus/chab5\_Aplus.gp}). The differentials
$\varphi_{i}^{*}\omega_{E_{i}}$ form a basis of $H^{0}(X_{M_{0}},\Omega^{1})$,
so for $z\in X(\Q)$, writing $[z-\infty_{+}]=\lambda D+\epsilon T$ with
$\lambda\in\Q$ and $\epsilon\in\Z/3$ (and $\Log T=0$),
\[
  (\mathcal{L}_{1}(z),\mathcal{L}_{2}(z))
   =\lambda\cdot(\mathcal{L}_{1}(D),\mathcal{L}_{2}(D)),\qquad \lambda\in\Q_{5},
\]
where $\mathcal{L}_{i}(D):=\mathcal{L}_{i}((-1,3))$. At $(-1,\pm3)$ this holds
with $\lambda=\pm1$ exactly, to $O(5^{44})$, and
$\rho(z):=\mathcal{L}_{2}(D)\mathcal{L}_{1}(z)-\mathcal{L}_{1}(D)\mathcal{L}_{2}(z)$
vanishes at all four known points. On each of the six residue disks of
$X(\Q_{5})$ --- $x\equiv1,4 \bmod 5$ with two signs of $y$, and the two disks at
infinity --- expand $\mathcal{L}_{i}$ by termwise integration of
$\varphi_{i}^{*}\omega_{E_{i}}$ and solve $\rho=0$. Because $E_{2}\cong E_{1}$
over $\Q_{25}$ the $\Q_{5}$-component of $\rho$ vanishes identically and the
condition is the $\sqrt2$-component; its Strassman indices are $1,1,2,2,1,1$,
giving exactly $8$ zeros in $X(\Q_{5})$:
\begin{itemize}[leftmargin=2.4em]
\item $\infty_{\pm}$ and $(-1,\pm3)$, with $\lambda=0,0,\pm1$;
\item $(-2\sqrt[3]{2},\pm12\sqrt6)$ with $\lambda=0$, a torsion class of order
      $3$ --- verified exactly over $\Q(\sqrt[6]{-2},\sqrt{-3})$
      (\texttt{aplus/tors12.gp}), so no rational point lies in its disk, that
      disk having a unique zero;
\item one pair $z_{1},\iota z_{1}$ in the disk $x\equiv1\pmod5$ with
      $\lambda\equiv\pm8316\pmod{5^{6}}$, so $\lambda\not\equiv0,\pm1\pmod{25}$.
\end{itemize}

\emph{(d) Sieve.} A rational point $z$ in the disk $x\equiv1$ would have
$\lambda(z)=n/m\in\Q$ with $5\nmid m$: the logarithms of $E_{i}(\Q_{25})$ have
valuation $\ge1$ (formal group, and $\#E_{i}(\FF_{25})=36$ is a $5$-unit) while
$\mathcal{L}_{i}(D)$ has valuation exactly $1$, so if $D=mG$ for a generator $G$
then $5\nmid m$. Hence for every prime $\qq$ of $M_{0}$ with
$5^{k}\mid \#E_{i}(\FF_{\qq})$ the $5$-primary parts satisfy
$\pi_{i}(\varphi_{i}(\bar z)-\varphi_{i}(\infty_{+}))=\lambda(z)\pi_{i}(\varphi_{i}(\bar D))$
in $E_{i}(\FF_{\qq})$ (the class $T$ has order $3$), with
$\lambda(z)\equiv\lambda(z_{1})\pmod{5^{6}}$ if $z$ lies in the disk of $z_{1}$.
The script \texttt{aplus/sieve5\_Aplus.gp} enumerates $A^{+}(\FF_{q})$ and
collects the residues $\lambda \bmod 5^{4}$ realised: at the three degree-$2$
primes above $q=149$ (where $\#E_{i}(\FF_{\qq})=22500=2^{2}3^{2}5^{4}$) $175$ of
the $625$ residues occur, and at a degree-$2$ prime above $q=599$ (where
$\#E_{i}=360000=2^{6}3^{2}5^{4}$) $275$ of $625$. The residues $0,\pm1$ of the
known points occur at both, while $\lambda(z_{1})\equiv191$ and
$\lambda(\iota z_{1})\equiv434 \pmod{625}$ occur at neither $q=599$. So the disk
$x\equiv1$ contains no rational point. The disks $x\equiv4$ have Strassman index
$2$ and contain the exhibited zeros $(-1,\pm3)$ and
$(-2\sqrt[3]{2},\pm12\sqrt6)$; the disks at infinity have index $1$ and contain
$\infty_{\pm}$.
\end{proof}

\begin{remark}
At a prime split in $\Q(\sqrt{-2})$ --- we first used $p=43$,
\texttt{aplus/linchab\_Aplus.gp} --- the same computation has $66$ zeros: the $4$
rational points, $12$ algebraic points carrying torsion classes ($x^{3}=2$,
$x^{3}=-16$, over $\Q(\sqrt[6]{-2},\sqrt{-3})$, of exact orders $3$ and $6$),
and $50$ others, of which the $43$-adic sieve removes all but $14$. The
survivors include the $\Q(\sqrt{-2})$-points $(2,\pm12\sqrt{-2})$ and
$(\pm2\sqrt{-2},\pm(16\pm20\sqrt{-2}))$, whose classes have $43$-adic logarithms
on the line of $J(\Q)$ because the CM action of $\Z[\sqrt{-2}]$ on $E_{i}$
commutes with the $43$-adic logarithms when $43$ splits in $\Q(\sqrt{-2})$:
Chabauty at a split prime sees $J(\Q(\sqrt{-2}))$ rather than $J(\Q)$. This is
why the inert prime $5$ is the right choice.
\end{remark}

% ==============================================================================
\section{The hard pair: automorphisms and N\'eron--Severi}\label{sec:structure}
% ==============================================================================

From here on $X$ denotes $B^{-}\colon y^{2}=-C_{3}(x)$; the argument for
$C^{+}\colon y^{2}=2C_{3}(x)$, its quadratic twist by $-2$, is identical with
the twisted data, and all numerical values quoted are those for $B^{-}$ unless
stated otherwise. $X$ has the rational points $P_{1}=(1,1)$, $P_{2}=(-1,15)$ and
their $\iota$-images; $J=\Jac(X)$.

\begin{proposition}\label{prop:aut}
\begin{enumerate}[label=(\roman*),leftmargin=2.4em]
\item $\Aut(X_{\Qbar})\cong\mathrm{GL}_{2}(\FF_{3})$, of order $48$; $X$ is a
      $\Q$-twist of the Bolza curve $y^{2}=x^{5}-x$.
\item The M\"obius maps permuting the six roots of $C_{3}$ form $\Sym_{4}$;
      nine of them are involutions, and the six that lift to involutions $\tau$
      of $X$ are $x\mapsto(-x+b)/(cx+1)$ with
      \[
        81b^{6}-540b^{5}+1494b^{4}-2240b^{3}+1980b^{2}-1008b+232=0,
      \]
      the lift being
      $(x,y)\mapsto\bigl(\mu(x),\,\kappa s\,y/(cx+1)^{3}\bigr)$, where
      $\kappa=1+bc$ and $s^{2}=\kappa$. (The other three involutory M\"obius
      maps, those with $b^{3}-18b^{2}+36b-20=0$, lift to automorphisms of order
      $4$.)
\item The M\"obius parts are defined over
      $K_{6}=\Q(\sqrt[3]{2},\sqrt{\sqrt[3]{2}-1})$ and the involutions over
      $L=K_{6}(\sqrt{\kappa})$, a field of degree $12$ with
      \[
        L=\Q[t]/(t^{12}+6t^{10}+3t^{8}-28t^{6}-33t^{4}+6t^{2}-3),
      \]
      of signature $(2,5)$, discriminant $2^{28}3^{13}$ and class number $1$, in
      which $2$ is totally ramified and $3=\pp_{1}^{3}\pp_{2}^{6}\pp_{3}^{3}$.
\item Over $L$, $X$ is bielliptic: with $x_{\pm}=(-1\pm s)/c$ the fixed points of
      $\mu$ and
      \[
        \mathcal{X}=\frac{x-x_{+}}{x-x_{-}},\qquad
        \mathcal{Y}=y(\mathcal{X}-1)^{3},
      \]
      the curve becomes
      $\mathcal{Y}^{2}=F(\mathcal{X})=A\mathcal{X}^{6}+B\mathcal{X}^{4}+C\mathcal{X}^{2}+D$
      with $A=f(x_{-})$, $D=f(x_{+})$, and
      $\tau(\mathcal{X},\mathcal{Y})=(-\mathcal{X},\mathcal{Y})$. The quotients
      \[
        E_{1}\colon Y'^{2}=X'^{3}+BX'^{2}+ACX'+A^{2}D,\qquad
        E_{2}\colon Y'^{2}=X'^{3}+CX'^{2}+BDX'+AD^{2},
      \]
      reached by $\varphi_{1}=(A\mathcal{X}^{2},A\mathcal{Y})$ and
      $\varphi_{2}=(D/\mathcal{X}^{2},D\mathcal{Y}/\mathcal{X}^{3})$, have
      $j=8000$ (CM by $\Z[\sqrt{-2}]$), admit global minimal models over $L$,
      have minimal discriminants supported on the prime above $2$ and one prime
      above $3$, and $E_{i}(L)_{\tors}\cong\Z/2$.
\end{enumerate}
\end{proposition}

\begin{proof}
Exact computation. \texttt{invol.gp} finds the $24$ M\"obius maps numerically,
and \texttt{invol4.gp}, \texttt{invol8.gp} identify them exactly by resultants:
the nine involutory maps have $b$ a root of
$(b^{3}-18b^{2}+36b-20)\cdot(81b^{6}-\cdots)\cdot(-C_{3}(b))$, the last factor
being spurious. One has $\lambda/(1+bc)^{3}=+1$ for the six genuine involutions
and $-1$ for the three lifts of order $4$. Fields, discriminants, class numbers,
minimal models and torsion: \texttt{invol8.gp}, \texttt{step2a.gp}. The
normalisation of the quotient maps follows \cite{BianchiPadurariu}.
\end{proof}

\begin{remark}
The natural quotient models are non-minimal at two degree-$1$ primes above
$397$, where $A$ and $D$ respectively have valuation $-6$. This is the source of
the constant $\Omega_{397}$ in Theorem~\ref{thm:identity}.
\end{remark}

\begin{proposition}[Galois action on $\NS$]\label{prop:NS}
Let $N$ be the $\Sym_{4}$-field obtained as the Galois closure of
$\Q(\sqrt[3]{2},\omega,\sqrt{-1-\sqrt[3]{2}})$, let $F=\Q(\sqrt{-2})$ and
$M=NF$, so $[M:\Q]=48$ and $\Gal(M/\Q)\cong\Sym_{4}\times\Z/2$. Then
\[
  \NS(J_{\Qbar})\otimes\Q\;\cong\;\mathbf{1}\oplus(\sgn_{F}\otimes\sgn\otimes\stdrep).
\]
In particular $\rho(J/\Q)=\rho(J/F)=1$, and complex conjugation acts with
eigenvalues $(1,1,1,-1)$: there is exactly one anti-invariant class, namely the
class $Z=\tau_{*}$ of an involution, defined over $L$.
\end{proposition}

\begin{proof}
Frobenius traces (\texttt{nschar2.gp}). Every prime inert in $F$ has trace $0$,
so $V_{\ell}(J)=\Ind_{F}^{\Q}\psi$; at split primes the Frobenius polynomial
factors over $F$ (constant terms $\pm\pi^{2},\pm\bar\pi^{2}$) but over $\Z$ only
when it is even. Hence $\End^{0}_{F}(J)=F$, and since the Rosati involution on
an imaginary quadratic field is complex conjugation, $\rho(J/F)=1$.

The trace of Frobenius on $\NS(J_{\Qbar})\otimes\Q_{\ell}$ --- the count of
root-of-unity eigenvalues among the $\alpha_{i}\alpha_{j}/p$, the $2$-dimensional
transcendental part contributing $0$ at inert primes --- depends only on the
splitting of $p$ in $F$ and in $N$:
\begin{center}\small
\begin{tabular}{@{}llcc@{}}
\toprule
class in $N$ & condition on $p$ & $F$-split & $F$-inert\\
\midrule
$3$-cycles        & $p\equiv1\ (3)$, $x^{3}-2$ irred.\ mod $p$              & $1$ & $1$\\
double transp.    & $p\equiv1\ (3)$, three roots, $(-1-r\mid p)=[-1,-1,1]$  & $0$ & $2$\\
identity          & $p\equiv1\ (3)$, three roots, all $(-1-r\mid p)=1$      & $4$ & $-2$\\
transpositions    & $p\equiv2\ (3)$, $(-1-r\mid p)=-1$                      & $0$ & $2$\\
$4$-cycles        & $p\equiv2\ (3)$, $(-1-r\mid p)=+1$                      & $2$ & $0$\\
\bottomrule
\end{tabular}
\end{center}
These are the values of the character of
$\mathbf{1}\oplus(\sgn_{F}\otimes\sgn\otimes\stdrep)$ for
$\Gal(M/\Q)\cong\Sym_{4}\times\Z/2$, and the observed class densities match the
orders of the conjugacy classes. The element (transposition, $\sigma_{F}$) is
complex conjugation and has trace $2$, so its eigenvalues are $(1,1,1,-1)$. The
stabiliser of the anti-invariant class has order $4$ in a group of order $48$,
giving a field of definition of degree $12$ --- which is the field $L$ of
Proposition~\ref{prop:aut}.
\end{proof}

\begin{remark}
Proposition~\ref{prop:NS} also records where local height contributions can
arise: at $3$ the stable model is smooth (potentially good reduction), and at
$2$ the stable model is a pair of supersingular elliptic curves.
\end{remark}

% ==============================================================================
\section{The quadratic Chabauty identity}\label{sec:identity}
% ==============================================================================

\subsection{The imported statements}

\begin{theorem}[Balakrishnan--Dogra {\cite[Prop.~2]{BDii}}, as quoted in {\cite{Dogra}}]
\label{thm:BD}
Let $X/\Q$ be a curve of genus $g$ with Jacobian $J$. The depth-$2$
Chabauty--Kim set $X(\Q_{p})_{2}$ is finite whenever
\[
  \rk J(\Q)\;<\;g-1+\rk\NS(J)+\rk\NS(J_{\Qbar})^{c=-1}.
\]
\end{theorem}

Here $2<2-1+1+1$, by Proposition~\ref{prop:NS}.

\begin{theorem}[Balakrishnan--Besser--Bianchi--M\"uller {\cite[Thm.~1.5]{BBBM}}]
\label{thm:BBBM}
Let $X$ be a bielliptic curve of genus $2$ over a number field $L$, with
quotient maps $\varphi_{1},\varphi_{2}$ to elliptic curves $E_{1},E_{2}$, and let
$\chi\colon \Ad_{L}^{\times}/L^{\times}\to\Q_{p}$ be a non-trivial continuous
idele class character. Then
\[
  \rho^{\chi}(z)=\sum_{\pp\mid p}\Bigl(h^{\chi}_{\pp}(\varphi_{1}(z_{\pp}))
   -h^{\chi}_{\pp}(\varphi_{2}(z_{\pp}))-2\chi_{\pp}(x(z_{\pp}))\Bigr)
   -\sum\alpha\, q_{1}(\varphi_{1}(z))+\sum\beta\, q_{2}(\varphi_{2}(z))
\]
maps $X(L)$ into an explicitly computable finite set $T^{\chi}$.
\end{theorem}

We apply Theorem~\ref{thm:BBBM} to $X_{L}$, restricted to
$X(\Q)\subset X(L)$, for one specific character. The role of the term
$\rk\NS(J_{\Qbar})^{c=-1}$ in Theorem~\ref{thm:BD} is the following selection
rule, which is what makes the restriction useful: for $z\in X(\Q)$ all
Galois-conjugate places contribute the same local terms, so the identity
survives on $\Q$-points only for pairs (class, character) lying in the same
Galois-isotypic component. The cyclotomic character pairs with the trivial
component and gives $0=0$; the usable classes are the
complex-conjugation-anti-invariant ones.

\subsection{The character}

\begin{proposition}\label{prop:character}
Let $p=499$, the smallest prime that splits completely in $M$, so that
$L\otimes\Q_{p}=\Q_{p}^{12}$ and every endomorphism is $p$-rational. Under the
twelve embeddings $\sigma_{i}\colon L\to\Q_{p}$ the involutions act on
$H^{0}(X,\Omega^{1})$, in the basis $dx/y,\,x\,dx/y$, by
\[
  A_{i}=s_{i}^{-1}\begin{pmatrix}-1&-c_{i}\\-b_{i}&1\end{pmatrix}.
\]
Because $H^{0}(\Omega^{1})$ is one of the two CM eigenspaces of Frobenius on
$H^{1}_{\mathrm{dR}}$, one has $\NS(J)\otimes\Q_{p}\cong M_{2}(\Q_{p})$ through
this action, and the $A_{i}$ span $\mathfrak{sl}_{2}$, the non-trivial
$3$-dimensional part. Let $R\subset\Q_{p}^{12}$ be the orthogonal complement of
$\ker\bigl((c_{i})\mapsto\sum_{i}c_{i}A_{i}\bigr)$; then $\dim R=3$ and $R$ is
the isotypic component of $\sgn_{F}\otimes\sgn\otimes\stdrep$ in the permutation
representation on the places above $p$. The six fundamental units of $L$ have
$p$-adic logarithms of rank $6$, meeting $R$ in a $2$-dimensional space. Hence
there is, up to scaling, exactly one idele class character $\chi_{Z}$ of $L$
whose local components at the places above $p$ lie in $R$:
\[
  \chi_{Z}=(c_{1},\dots,c_{12}),\quad c_{1}=1,
\]
extended to all places by $h_{L}=1$ via
$\chi_{w}(\pi_{w})=\sum_{i}c_{i}\log_{p}\sigma_{i}(\alpha_{w})$ for a generator
$\alpha_{w}$ of $w$. Moreover $\chi_{Z}$ vanishes at the uniformisers of every
prime of $L$ above $2$ and above $3$.
\end{proposition}

\begin{proof}
\texttt{step1\_character.gp}: the kernel has dimension $9$, so $\dim R=3$;
$R\perp(1,\dots,1)$; the unit logarithms have rank $6$ and
$\dim(R\cap\operatorname{span}(\text{unit logs}))=2$; and the values at the four
uniformisers above $2$ and $3$ are $O(499^{30})$.
\end{proof}

\begin{remark}
The vanishing at $2$ and $3$ is what makes the identity below independent of the
reduction types there: no analysis of the bad fibres is required. Whether it has
a structural explanation is one of the questions the author would put to an
expert.
\end{remark}

\subsection{The identity}

Set $D_{1}=[P_{1}-\iota P_{1}]$, $D_{2}=[P_{2}-\iota P_{2}]\in J(\Q)$ and
$Q_{i}=\varphi_{1}(P_{i})\in E_{1}(L)$, $Q'_{i}=\varphi_{2}(P_{i})\in E_{2}(L)$.
The $499$-adic logarithms of $Q_{1},Q_{2}$ at two places are independent, so
$D_{1},D_{2}$ span $J(\Q)\otimes\Q$; the rank is exactly $2$ by
\S\ref{sec:rank}.

\begin{theorem}\label{thm:identity}
For $z\in X(\Q)$ write $[z-\iota z]=a_{1}D_{1}+a_{2}D_{2}$ in $J(\Q)\otimes\Q$
and set $a(z)=(a_{1},a_{2})^{\mathsf T}$, and let
\[
  G_{ij}=\langle Q_{i},Q_{j}\rangle_{\chi}-\langle Q'_{i},Q'_{j}\rangle_{\chi}
\]
be the difference of the $\chi_{Z}$-height pairings on $E_{1}(L)$ and
$E_{2}(L)$. Then for every $z\in X(\Q)$,
\begin{equation}\label{eq:identity}
  \rho(z):=\sum_{i=1}^{12}c_{i}\bigl[\lambda_{i}(\varphi_{1}z)-\lambda_{i}(\varphi_{2}z)\bigr]
  \;-\;2\sum_{i=1}^{12}c_{i}\log_{p}\sigma_{i}(\mathcal{X}(z))
  \;-\;a(z)^{\mathsf T}Ga(z)\;=\;-\Omega_{397},
\end{equation}
where $\lambda_{i}$ is the canonical Mazur--Tate local height at the $i$-th place
above $p$ on the $b$-model of the relevant quotient, and $\Omega_{397}$ is the
constant $2\chi'_{w_{1}}(\pi)-2\chi'_{w_{2}}(\pi)$ attached to the two degree-$1$
primes $w_{1},w_{2}$ of $L$ above $397$.
\end{theorem}

\begin{proof}
Since $\varphi_{1}(z)=a_{1}Q_{1}+a_{2}Q_{2}+T$ with $T\in E_{1}(L)[2]$, and
heights kill torsion, $h_{\chi}(\varphi_{1}z)-h_{\chi}(\varphi_{2}z)=a^{\mathsf T}Ga$.
Decompose $h_{\chi}$ into local terms: at the twelve places above $p$,
$c_{w}\lambda_{w}(Q)+d_{w}m_{w}(Q)$ with $m_{w}=\max(0,-v_{w}(x(Q)))$; at a
place $w\nmid p$ of good reduction of the minimal model, $\chi'_{w}(\pi_{w})m_{w}(Q)$;
and at the places above $2$ and $3$, nothing, by Proposition~\ref{prop:character}.

On the natural quotient models $x(\varphi_{1}z)=A\mathcal{X}^{2}$ and
$x(\varphi_{2}z)=D/\mathcal{X}^{2}$, so at every place $w$ not above $2,3,397$ ---
where $A,D$ are units and the models are minimal ---
\begin{equation}\label{eq:mdiff}
  m_{w}(\varphi_{1}z)-m_{w}(\varphi_{2}z)=-2v_{w}(\mathcal{X}(z)),
\end{equation}
and the same holds at $w\mid p$. Summing $\chi'_{w}(\pi_{w})\cdot(-2v_{w}(\mathcal{X}(z)))$
over $w\nmid p$, and using that $\chi_{Z}$ kills the global element
$\mathcal{X}(z)\in L^{\times}$, produces
\[
  -2\sum_{i}c_{i}\log_{p}\sigma_{i}(\mathcal{X}(z))
  \;+\;2\sum_{i}d_{i}v_{p}(\sigma_{i}\mathcal{X}(z)),
\]
and the $d_{i}$-terms cancel against the terms
$d_{i}(m_{i}(\varphi_{1}z)-m_{i}(\varphi_{2}z))$ at $p$, by \eqref{eq:mdiff}
again. What remains is the contribution of the two primes above $397$, where the
natural models are not minimal; this is Lemma~\ref{lem:397} below, which shows
it to be the constant $\Omega_{397}$.
\end{proof}

\begin{lemma}\label{lem:397}
Let $w$ be one of the two degree-$1$ primes of $L$ above $397$, with
$v_{w}(A)=-6$ at the first and $v_{w}(D)=-6$ at the second, and let $u_{1},u_{2}$
be the scalings to the minimal models, of valuations $v_{w}(u_{1})=-2$ and
$v_{w}(u_{2})=-1$ (discriminant exponents $-24$ and $-12$). Then, computing
$m_{w}$ on the minimal models,
\[
  m_{w}(\varphi_{1}z)-m_{w}(\varphi_{2}z)+2v_{w}(\mathcal{X}(z))=+2
\]
at the first prime, for every $397$-adic point $z$, and $-2$ at the second.
Hence the correction in \eqref{eq:identity} is the constant
$\Omega_{397}=2\chi'_{w_{1}}(\pi)-2\chi'_{w_{2}}(\pi)$.
\end{lemma}

\begin{proof}
Case analysis on $v:=v_{w}(\mathcal{X}(z))$, using $v_{w}(x_{-})=-1$ and
$x_{+}$ integral. For $v\le0$: $m_{w}(\varphi_{1})=2-2v$ and
$m_{w}(\varphi_{2})=0$. For $v=1$: both vanish. For $v\ge2$:
$m_{w}(\varphi_{1})=0$ and $m_{w}(\varphi_{2})=2v-2$. In each case the stated
combination equals $2$; the value $-2$ at the second prime follows by the
conjugation symmetry exchanging the roles of $A$ and $D$.
\end{proof}

\begin{remark}
The constant is computed from a generator of the two primes and from
$\mathcal{X}(P_{1})$. Numerically $\rho(P_{k})+\Omega_{397}=O(499^{19})$ at both
known points with the sign $-2\log\mathcal{X}$ of \eqref{eq:identity}, and this
fails with the opposite sign (\texttt{step2b.gp}); so the numerics pin the
normalisation independently of the derivation.
\end{remark}

\subsection{Computing the heights}

Throughout we use the $\sigma$-function description of the local height at $p$,
which replaces a double Coleman integral by an evaluation of the Mazur--Tate
$\sigma$-function \cite{MazurTate} and is what makes the computation below
practical; we follow \cite{Bianchi} for that description and \cite{BBBM} for its
extension to a number field and a non-cyclotomic character.

The heights were computed as follows (\texttt{step2b.gp}). Each
$E_{k}\otimes_{\sigma_{i}}\Q_{p}$ is replaced by an integer model congruent to it
modulo $499^{30}$, so that PARI's \texttt{ellpadics2} supplies the $\sigma$-constant;
$\sigma$ itself comes from the formal group. With $m=\#E(\FF_{p})=514$, the
division polynomial $\psi_{m}(Q)$ is evaluated by the classical recurrences
rather than by expanding a polynomial of degree ${\approx}130{,}000$, calibrated
by
\[
  \log\sigma(kQ_{0})-k^{2}\log\sigma(Q_{0})=\log\psi_{k}(Q_{0})
  \qquad (k=2,3,4)
\]
to $O(499^{19})$ --- which also fixes PARI's division-polynomial parity
convention and the $a_{1}=a_{3}=0$ requirement of the $\sigma$-identities ---
and then
\[
  \lambda(Q)=-\frac{2}{m^{2}}\bigl(\log\sigma(mQ)-\log\psi_{m}(Q)\bigr).
\]

\emph{Validation.} The $\chi_{Z}$-height is a quadratic form: on both quotients
\[
  h(Q_{1}+Q_{2})+h(Q_{1}-Q_{2})-2h(Q_{1})-2h(Q_{2})=O(499^{19}),
  \qquad h(2Q_{1})=4h(Q_{1}).
\]
The relative sign between the $p$-adic and away-from-$p$ terms was wrong in a
first version, and this test is what caught it.

% ==============================================================================
\section{Zeros and the Mordell--Weil sieve}\label{sec:sieve}
% ==============================================================================

\begin{proposition}\label{prop:zeros}
$\rho$ extends to a locally analytic function on $X(\Q_{p})$. On the disks where
some $\sigma_{i}(\mathcal{X}(z))$ tends to $0$ or $\infty$, the singular pair of
terms is replaced by its analytic form $2c_{i}[\log(s/\mathcal{X})+v(s)]$ with
$s$ the formal parameter, as in \cite[Rem.~2.5]{BianchiPadurariu}; the
logarithmic singularities of $\lambda_{i}$ and of $\log\mathcal{X}$ cancel.

At $p=499$, $X(\Q_{499})$ has $528$ residue disks and no Weierstrass disk
($C_{3}$ has no root mod $499$). Expanded to $O(T^{14})$ with $14$-digit
coefficients, $\rho$ has Strassman index $2$ on every disk, with coefficient
valuations $[\,\cdot\,,1,1,2,3,\dots]$ --- the quadratic term dominates because
every logarithm lies in $499\Z_{499}$ and the inverse log-matrix has valuation
$-1$. The $\Z_{499}$-roots of the truncated series number $524$ for $B^{-}$
($262$ disks with two roots, $266$ with none) and $504$ for $C^{+}$.
Recomputation with $20$ digits and $O(T^{20})$ reproduces the roots to $8$
digits and the counts, the valuations continuing linearly to $18$, so the
truncation loses nothing. Finally $\rho$ vanishes to $O(499^{13})$ at the known
points.
\end{proposition}

\begin{proof}
The expansion and the root-finding are \texttt{step3\_setup.gp},
\texttt{step3\_disk.gp} and \texttt{step3\_run.gp}, aggregated by
\texttt{step3\_aggregate.py}; the higher-precision rerun is
\texttt{step3\_\allowbreak precision\_check.gp}. The per-disk output is in
\texttt{logs/\allowbreak step3\_disks.txt}, at $0.74$ s per disk, run in four
parallel chunks.
\end{proof}

\begin{proposition}[The sieve]\label{prop:sieve}
Every point of $J(\Q_{499})$ has a $499$-integral coefficient vector $a$: all
logarithms lie in $499\Z_{499}$, and the log-matrix of $Q_{1},Q_{2}$ has
determinant of valuation $2$. Hence
$[J(\Q):\langle D_{1},D_{2}\rangle+J(\Q)_{\tors}]$ is a $499$-unit, and for a
degree-$1$ prime $\qq$ of $L$ above $q$ with
$499\mid\#E_{1}(\FF_{\qq})=\#E_{2}(\FF_{\qq})$, projecting to the $499$-parts
$\xi\colon E_{k}(\FF_{\qq})\to\Z/499$, a rational point $z$ must satisfy
\[
  \xi\bigl(a(z)\cdot(\bar Q,\bar Q')\bigr)
   \in\bigl\{\xi(\varphi_{1}(w),\varphi_{2}(w)) : w\in X(\FF_{q})\bigr\}
   \subset(\Z/499)^{2},
\]
with no torsion and no index ambiguity, both $2$-torsion translates and unit
indices vanishing mod $499$.

Such primes $q$ are: four above $997$ (with $\#E=998$), four above $12973$, two
above $47903$, two above $59879$, and twelve above $88339$. The first prime above
$997$ eliminates $518$ of the $524$ zeros on $B^{-}$ (and $494$ of $504$ on
$C^{+}$); the primes above $12973$ eliminate all the rest but four. On each
curve the four survivors are the four known points: $(1,\pm1)$ and $(-1,\pm15)$
with $a=(\pm1,0),(0,\pm1)$ on $B^{-}$, and $(2,\pm12)$, $(6/5,\pm172/125)$
likewise on $C^{+}$.
\end{proposition}

\begin{proof}
\texttt{step4\_scan.gp}, \texttt{step4\_sieve.gp}; \texttt{logs/step4\_sieve.log}.
\end{proof}

\subsection{An independent second prime}

The whole of \S\ref{sec:identity}--\S\ref{sec:sieve} was repeated at $p=1459$,
the next prime that splits completely in $L$ and at which $C_{3}$ has no root
(\texttt{second\_prime/\allowbreak make\_p1459.py}, which regenerates the
pipeline from the $p=499$ scripts; logs in \texttt{logs/p1459/}). The structural
facts recur: the
character is again unique in its isotypic component and vanishes at the
uniformisers above $2$ and $3$; the $\chi_{Z}$-heights satisfy the parallelogram
law to $O(1459^{19})$; $\rho(P_{1})=\rho(P_{2})=-\Omega_{397}$ to
$O(1459^{19})$ on both curves; and $\#E_{k}(\FF_{\pp})=1462$ at all twelve
places, so the reduction is ordinary. Each curve has $1464$ residue disks, all of
Strassman index $2$ except the two disks at infinity (index $1$), and $\rho$ has
$1546$ zeros on $B^{-}$ and $1426$ on $C^{+}$. The sieve primes are the degree-$1$
primes of $L$ above $q=2917,\,26261,\,96293$, all $\equiv-1\pmod{1459}$ with
$\#E_{1}(\FF_{\qq})=\#E_{2}(\FF_{\qq})=q+1$; the first prime above $2917$
removes $1538$ of the $1546$ zeros on $B^{-}$, and $26261$ removes all the rest
but four, which are the known points.

On $C^{+}$ the same primes leave the four known points together with the two
zeros of each of the two disks $x\equiv103\pmod{1459}$. On those disks the
logarithms at places $1$ and $2$ lose all precision, and the $14$-digit run had
returned coefficient vectors with no significant digits. Recomputing them with
$22$ working digits (\texttt{logs/p1459/Cplus/disk103.log} and the two follow-up
logs) gives the same two zeros to five digits and ordinary coefficient vectors
$a\equiv(243,104),(1453,735),(1216,1355),(6,724)\pmod{1459}$, which the prime
above $2917$ removes. So at $p=1459$ too, $B^{-}(\Q)$ and $C^{+}(\Q)$ consist of
exactly the four known points each.

\begin{remark}
This is the one place where the pipeline required a per-disk precision
intervention. A systematic per-disk precision check on $a(z)$, and not only on
$\rho$, would be the right way to handle it; the present runs check it by hand
on the one disk where it was needed.
\end{remark}

\begin{corollary}\label{cor:hardpair}
Granting $\rk J(\Q)=2$ for both Jacobians (\S\ref{sec:rank}),
\[
  B^{-}(\Q)=\{(1,\pm1),\,(-1,\pm15)\},\qquad
  C^{+}(\Q)=\{(2,\pm12),\,(6/5,\pm172/125)\}.
\]
\end{corollary}

% ==============================================================================
\section{The rank bounds}\label{sec:rank}
% ==============================================================================

\begin{proposition}\label{prop:rank}
For $X=B^{-}$ and for $X=C^{+}$: $J(\Q)[2]=0$ and $\dim_{\FF_{2}}\Sel^{2}(J/\Q)=2$.
Hence $\rk J(\Q)\le2$, so $\rk J(\Q)=2$ and $\Sha(J/\Q)[2]=0$.
\end{proposition}

\begin{proof}
Write $f$ for the sextic ($-C_{3}$ for $B^{-}$, $2C_{3}$ for $C^{+}$), let
$L_{6}=\Q[x]/(f)$ and let $\vartheta$ be the image of $x$. In both cases $L_{6}$
has class number $1$, signature $(2,2)$ and Galois group $2\wr \Sym_{3}$ of order
$48$, and no quadratic subfield (\texttt{nfsubfields}).

\emph{(a) Torsion.} $J[2](\Qbar)$ is the group of even-cardinality subsets of the
roots of $f$ modulo complementation, so a $\Q$-rational $2$-torsion point is a
Galois-stable such subset, that is, a union of irreducible factors. Since $f$ is
irreducible, $J(\Q)[2]=0$.

\emph{(b) The fake Selmer group.} Let
$H=\ker\bigl(\Nm\colon L_{6}^{\times}/L_{6}^{\times2}\Q^{\times}\to\Q^{\times}/\Q^{\times2}\bigr)$.
The map $D=\sum n_{P}P\mapsto\prod(x_{P}-\vartheta)^{n_{P}}$ induces a
homomorphism $x-T\colon J(\Q)\to H$, and likewise over every completion, and the
fake $2$-Selmer group is
\[
  \Sel^{2}_{\fake}(J)=\{\delta\in H:\delta_{v}\in(x-T)(J(\Q_{v}))\ \text{for all }v\}
\]
\cite[\S\S12--13]{PoonenSchaefer}, \cite[\S5]{Stoll2001}; see also
\cite{Creutz}. It is a quotient of the
true Selmer group: there is an exact sequence
$\mu_{2}\to\Sel^{2}(J)\to\Sel^{2}_{\fake}(J)\to0$
(\cite[\S5]{StollvanLuijk}, with \cite[Thm.~13.2]{PoonenSchaefer} for the
injectivity of the first map), whence
\[
  \dim\Sel^{2}(J)\le\dim\Sel^{2}_{\fake}(J)+1 .
\]
(In fact this is an equality here, the first map being injective because $f$ has
no odd-degree factor over $\Q$ and is not a constant times a product of two
conjugate cubics over a quadratic field, and the image of $\mu_{2}$ being the
class of $[P-\iota P]$ for a rational point $P$, which is a Selmer element. Only
the inequality is used.)

\emph{(c) Local images.} For a finite place $v$, let $m_{v}$ be the number of
irreducible factors of $f$ over $\Q_{v}$ and $o_{v}\in\{0,1\}$ indicate whether
one has odd degree. Then $\dim J(\Q_{v})[2]=m_{v}-1-o_{v}$, by the same count as
in (a), and $\dim J(\Q_{v})/2J(\Q_{v})=d_{v}:=\dim J(\Q_{v})[2]+2\cdot[v=2]$.
The kernel of $x-T$ on $J(\Q_{v})/2J(\Q_{v})$ is generated by the class $c_{v}$
of $[P-\iota P]$ for $P\in X(\Q_{v})$, which does not depend on $P$ (the
difference of two such classes is $2[P-P']$); and $c_{v}=0$ exactly when the
canonical class is twice a $\Q_{v}$-rational degree-$1$ class, i.e.\ when one of
the $16$ theta characteristics --- the six $[w_{i}]$ and the ten
$[w_{i}+w_{j}-w_{k}]$, one for each partition of the Weierstrass points into two
triples --- is Galois-invariant, i.e.\ when $o_{v}=1$ or all local factors of $f$
share a quadratic subfield $\Q_{v}(\sqrt{d})$ (the two orbits of
$\Gal(\Qbar_{v}/\Q_{v}(\sqrt d))$ on the roots of each factor then assembling
into a stable partition into triples). Writing $\epsilon_{v}=1$ in the remaining
case, the image $(x-T)(J(\Q_{v}))$ has dimension exactly $d_{v}-\epsilon_{v}$;
the computation generates it by $\Q_{v}$-points and conjugate pairs of quadratic
points and stops when that dimension is reached, so the local conditions imposed
are exactly the fake Selmer conditions. At the real place the image is given by
sign vectors \cite[\S5]{BruinStoll}, and is trivial because $f$ has two real
roots. At $p\nmid2\disc(f)$ the image is the unramified subgroup, so
$\Sel^{2}_{\fake}(J)\subseteq H(S):=\{\delta\in H: v_{\pp}(\delta)\text{ even for }\pp\nmid S\}$
for any $S\supseteq\{p\mid 2\disc(f)\}=\{2,3\}$; we used $S=\{2,3,17\}$, the
prime $17$ dividing the leading coefficient, which is harmless.

\emph{(d) Computation} (\texttt{sel2/pipeline.gp}; logs in
\texttt{logs/sel2/\allowbreak descent\_Bminus.log} and
\texttt{descent\_Cplus.log}). $L_{6}$ is
taken in a reduced model (\texttt{polredbest}; this matters --- in the model
$c^{5}f(y/c)$ PARI's local Hilbert symbols at the primes above $2$ are
unreliable). $H(S)$ is computed from the $S$-units, the class number being $1$,
as the norm-square subgroup of $L_{6}(S,2)$ modulo the image of $\Q(S,2)$, with
an assertion that the exponent vectors reconstruct $-1$ and the primes of $S$.
Local square classes at $\PPr\mid p$ are coordinates with respect to a basis of
$L_{6,\PPr}^{\times}/L_{6,\PPr}^{\times2}$ chosen by the rank of the Hilbert
pairing matrix (\texttt{nfhilbert}), and the local symbols are checked against
the product formula on $60$ random pairs.

For $B^{-}$: $\dim H(S)=2$. At $p=3$, $f$ is irreducible over $\Q_{3}$ and $3$ is
a square in $L_{6}\otimes\Q_{3}$, so $d=0$, $\epsilon=0$ and the image is $0$. At
$p=17$ the factors have degrees $1,1,4$, so $d=1$, $\epsilon=0$, and the image
has dimension $1$, generated by $\Q_{17}$-points. At $p=2$, $f$ is irreducible
over $\Q_{2}$ and no non-square of $\Q_{2}$ becomes a square in
$L_{6}\otimes\Q_{2}$, so $d=2$, $\epsilon=1$, and the image has dimension $1$,
generated by conjugate pairs of quadratic points ($72289$ pairs tested; the
$\Q_{2}$-points give nothing). The subgroup of $H(S)$ satisfying the three
conditions has dimension $1$, so $\dim\Sel^{2}_{\fake}(J)=1$ and
$\dim\Sel^{2}(J)\le2$. Since $\rk J(\Q)\ge2$, by the independent points of
\S\ref{sec:identity}, and $J(\Q)[2]=0$, we get $\dim\Sel^{2}(J)=2$ and
$\Sha[2]=0$.

For $C^{+}$: $\dim H(S)=2$ and the local data are the same (image of dimension
$1$ at $2$, from $151$ quadratic pairs; $0$ at $3$; $1$ at $17$), so
$\dim\Sel^{2}_{\fake}=1$ and $\dim\Sel^{2}=2$.
\end{proof}

\subsection{Validation of the implementation}

The pipeline was run on $24$ curves of the genus-$2$ database of Booker et
al.\ \cite{BSSVY} having rational points and irreducible sextic models
(\texttt{sel2/validation\_curve\_*.gp}; scorecard in
\texttt{logs/sel2/validation\_scorecard.txt}). For a curve with a rational point
and no $2$-torsion one has $\dim\Sel^{2}(J)=\rk J(\Q)+\dim\Sha[2]$ with
$\dim\Sha[2]$ even \cite{PoonenStoll}, so $\dim\Sel^{2}\ge\rk J(\Q)$ with
equality exactly when $\Sha[2]=0$. On all $24$ curves the computed
$\dim\Sel^{2}$ equals the database's analytic rank --- including curves with four
real roots, with class groups $\Z/2$ and $\Z/3$, with a quadratic subfield, and
of ranks $1$, $2$ and $3$.

Two curves with leading coefficient $4$, a $2$-adic square, needed
$\Q_{2}$-points near infinity, with $x$ of denominator up to $2^{9}$, to
complete the $2$-adic local image; the pipeline searches that chart when the
standard search falls short, and flags any image that remains incomplete, in
which case its output is only a lower bound.

Three earlier bugs --- the exponent ordering of \texttt{bnfissunit}, the Hilbert
symbols in the non-reduced model, and an empty local image at a prime with
trivial target --- had produced $\rk+1$ on ten of the validation curves, which the
parity of $\dim\Sha[2]$ rules out; that is how they were found.

% ==============================================================================
\section{Proof of the Theorem}\label{sec:conclusion}
% ==============================================================================

\begin{proof}[Proof of Theorem~\ref{thm:main}]
By Lemmas~\ref{lem:A1}--\ref{lem:A3}, every primitive solution of
\eqref{eq:main} gives a rational point on one of the six curves \eqref{eq:Apm},
\eqref{eq:BCpm}. By Proposition~\ref{prop:B1}, three of them have no rational
point; by Proposition~\ref{prop:B2}, $A^{+}(\Q)=\{\infty_{\pm},(-1,\pm3)\}$,
which contributes only $(\pm1,0,-1)$. By Corollary~\ref{cor:hardpair} together
with Proposition~\ref{prop:rank}, $B^{-}(\Q)$ and $C^{+}(\Q)$ consist of the
four known points each; transporting them back through Lemma~\ref{lem:A3},
the points $(1,\pm1)$ give $(\pm1,-1,1)$, the points $(-1,\pm15)$ and
$(2,\pm12)$ give an imprimitive solution, and $(6/5,\pm172/125)$ gives
$(\pm43,-203,237)$. That exhausts the primitive solutions, and each of the three
listed is checked directly to satisfy \eqref{eq:main}.
\end{proof}

The result is subject to the imported statements listed in
\S\ref{subsec:rests}: the theorems of Balakrishnan--Dogra and of
Balakrishnan--Besser--Bianchi--M\"uller, the standard theory of Coleman
integration used in Proposition~\ref{prop:B2} and in \S\ref{sec:identity}, and
the theory of $p$-adic heights on elliptic curves over number fields. The
$\sigma$-function local heights used here were validated against PARI's
canonical heights over $\Q$ in the companion implementation
(\texttt{supporting/qc\_sigma.py}). The $p$-adic computations carry PARI's
rigorous precision tracking; the only truncation is the series tail on each
residue disk, controlled by the observed linear growth of coefficient valuations
and confirmed at higher precision.

% ==============================================================================
\section{Code and data}\label{sec:code}
% ==============================================================================

All scripts and the logs of every run cited above are public at
\url{https://github.com/marchfederico/fermat433}: the code under the GNU
General Public License, version 3 or later, and this paper and the other
written documents under CC BY 4.0. The scripts in \texttt{exploration/}, which
nothing here depends on, are Denis Simon's and carry his copyright.
Everything runs on PARI/GP
2.17.4; the Python scripts need only \texttt{sympy}. No Magma or Sage is used.
The long runs are the residue-disk expansions (\texttt{step3\_run.gp} in four
chunks, about $10$ minutes at $p=499$ and about $40$ minutes at $p=1459$) and
the $2$-adic point searches of \texttt{sel2/pipeline.gp}, up to about $1.5$
hours per curve.

\begin{table}[ht]
\caption{Where each statement is established.}\label{tab:scripts}
\footnotesize
\renewcommand{\arraystretch}{1.15}
\begin{tabular}{@{}>{\raggedright\arraybackslash}p{0.19\textwidth}
                  >{\raggedright\arraybackslash}p{0.44\textwidth}
                  >{\raggedright\arraybackslash}p{0.29\textwidth}@{}}
\toprule
Statement & Scripts & Logs \\
\midrule
Lemmas~\ref{lem:A1}--\ref{lem:A2} &
  \texttt{descent\_explore.py}, \texttt{diagonal.py} & --- \\
Lemma~\ref{lem:A3} &
  \texttt{etale.py}, \texttt{twistset.py} & --- \\
Prop.~\ref{prop:B1} &
  \texttt{sel2/certify.gp}, \texttt{sel2/fakeselmerset.gp} &
  \texttt{logs/sel2/certify.log} \\
Prop.~\ref{prop:B2} &
  \texttt{aplus/bi\_Aplus.gp}, \texttt{aplus/chab5\_Aplus.gp},
  \texttt{aplus/sieve5\_Aplus.gp}, \texttt{aplus/tors12.gp} &
  \texttt{logs/aplus/} \\
Prop.~\ref{prop:aut} &
  \texttt{invol.gp}, \texttt{invol4.gp}, \texttt{invol8.gp}, \texttt{step2a.gp} & --- \\
Prop.~\ref{prop:NS} &
  \texttt{nschar2.gp}, \texttt{splitF.gp} & --- \\
Prop.~\ref{prop:character} &
  \texttt{step1\_character.gp} & \texttt{logs/step1\_data.txt} \\
Thm.~\ref{thm:identity} &
  \texttt{step2a.gp}, \texttt{step2b.gp} & \texttt{logs/step2b\_data.txt} \\
Prop.~\ref{prop:zeros} &
  \texttt{step3\_*.gp}, \texttt{step3\_aggregate.py} &
  \texttt{logs/step3\_disks.txt} \\
Prop.~\ref{prop:sieve} &
  \texttt{step4\_scan.gp}, \texttt{step4\_sieve.gp} &
  \texttt{logs/step4\_sieve.log} \\
Second prime &
  \texttt{second\_prime/make\_p1459.py} & \texttt{logs/p1459/} \\
$C^{+}$ throughout \S\S\ref{sec:structure}--\ref{sec:sieve} &
  \texttt{twist/} & \texttt{logs/twist/} \\
Prop.~\ref{prop:rank} &
  \texttt{sel2/pipeline.gp} & \texttt{logs/sel2/} \\
Validation &
  \texttt{sel2/validation\_\allowbreak curve\_*.gp} &
  \texttt{logs/sel2/\allowbreak validation\_\allowbreak scorecard.txt} \\
\bottomrule
\end{tabular}
\end{table}

\section*{Acknowledgements}

The collaboration with Claude (Anthropic) in which this work was done is
described in \S\ref{subsec:status}. The independent Magma computations of
G.~Villines \cite{Villines} agree with the present ones wherever the two
overlap, and that reduction is the starting point for \S\ref{sec:descent}.

% ================================================================== biblio ===
\begin{thebibliography}{BBBM}

\bibitem[BBBM]{BBBM}
J.~S. Balakrishnan, A.~Besser, F.~Bianchi, J.~S. M\"uller,
\emph{Explicit quadratic Chabauty over number fields},
Israel J. Math. \textbf{243} (2021), no.~1, 185--232. arXiv:1910.04653.

\bibitem[BD]{BDii}
J.~S. Balakrishnan, N.~Dogra,
\emph{Quadratic Chabauty and rational points II: Generalised height functions on
Selmer varieties},
Int. Math. Res. Not. IMRN \textbf{2021}, no.~15, 11923--12008.
arXiv:1705.00401.

\bibitem[Bi]{Bianchi}
F.~Bianchi,
\emph{Quadratic Chabauty for (bi)elliptic curves and Kim's conjecture},
Algebra Number Theory \textbf{14} (2020), no.~9, 2369--2416. arXiv:1904.04622.

\bibitem[BP]{BianchiPadurariu}
F.~Bianchi, O.~Padurariu,
\emph{Rational points on rank 2 genus 2 bielliptic curves in the LMFDB}, in:
LuCaNT: LMFDB, computation, and number theory, Contemp. Math. \textbf{796},
Amer. Math. Soc., Providence, RI, 2024, pp.~215--242. arXiv:2212.11635.

\bibitem[BSSVY]{BSSVY}
A.~R. Booker, J.~Sijsling, A.~V. Sutherland, J.~Voight, D.~Yasaki,
\emph{A database of genus-2 curves over the rational numbers},
LMS J. Comput. Math. \textbf{19} (2016), 235--254. arXiv:1602.03715.

\bibitem[BS]{BruinStoll}
N.~Bruin, M.~Stoll,
\emph{Two-cover descent on hyperelliptic curves},
Math. Comp. \textbf{78} (2009), no.~268, 2347--2370. arXiv:0803.2052.

\bibitem[Cr]{Creutz}
B.~Creutz,
\emph{Explicit descent in the Picard group of a cyclic cover of the projective
line}, in: ANTS X --- Proceedings of the Tenth Algorithmic Number Theory
Symposium, Open Book Ser.~\textbf{1}, Math. Sci. Publ., 2013.
arXiv:1204.5803.

\bibitem[Do]{Dogra}
N.~Dogra,
\emph{Unlikely intersections and the Chabauty--Kim method over number fields},
Math. Ann. \textbf{389} (2024), 1--62. arXiv:1903.05032.

\bibitem[MT]{MazurTate}
B.~Mazur, J.~Tate,
\emph{The $p$-adic sigma function},
Duke Math. J. \textbf{62} (1991), no.~3, 663--688.

\bibitem[PS]{PoonenSchaefer}
B.~Poonen, E.~Schaefer,
\emph{Explicit descent for Jacobians of cyclic covers of the projective line},
J. reine angew. Math. \textbf{488} (1997), 141--188.

\bibitem[PSt]{PoonenStoll}
B.~Poonen, M.~Stoll,
\emph{The Cassels--Tate pairing on polarized abelian varieties},
Ann. of Math. (2) \textbf{150} (1999), no.~3, 1109--1149.

\bibitem[RG]{RatcliffeGrechuk}
A.~Ratcliffe, B.~Grechuk,
\emph{Generalized Fermat equation: a survey of solved cases},
Expo. Math. (2025), article 125688. arXiv:2412.11933.

\bibitem[St]{Stoll2001}
M.~Stoll,
\emph{Implementing 2-descent for Jacobians of hyperelliptic curves},
Acta Arith. \textbf{98} (2001), no.~3, 245--277.

\bibitem[SvL]{StollvanLuijk}
M.~Stoll, R.~van Luijk,
\emph{Explicit Selmer groups for cyclic covers of $\PP^{1}$},
Acta Arith. \textbf{159} (2013), no.~2, 133--148. arXiv:1108.3364.

\bibitem[Vi]{Villines}
G.~Villines,
\emph{The Diophantine equation $x^{4}+2y^{3}+z^{3}=0$: reduction and partial
resolution}, Zenodo, \texttt{10.5281/zenodo.20672504} (June 2026).

\end{thebibliography}

\end{document}
